\paragraph{潜在生産水準目標（POLT）}
\label{sec:chapter5_beta_shock_policies_polt}
式
\eqref{form:chapter3_policies_monetary_home_polt_i_H_eq}および
\eqref{form:chapter3_policies_monetary_home_polt_gamma_H_eq}
により記述されたとおり、POLTの式は以下のようであった。
\begin{align}
\label{form:chapter5_beta_shock_policies_polt_i_H_eq}
    i_t^H&=i_{ss}^H+\phi_{gap}^H(\log(y_t^H)-\log(y_t^{H,pot}))+\phi_{level}^H\gamma_t^H \\
\label{form:chapter5_beta_shock_policies_polt_gamma_H_eq}
    \gamma_t^H&=\gamma_{t-1}^H+(\log(y_{t-1}^H)-\log(y_{t-1}^{H,pot}))
\end{align}
ここで
\(\phi_{gap}^H,\phi_{level}^H\)は反応係数、
\(y_t^H\)は実質産出量、
\(y_t^{H,pot}\)は潜在産出量、
\(\gamma_t^H\)は過去の乖離の累積である。
\begin{enumerate}
    \item
    \label{itm:chapter5_beta_shock_policies_polt_c_H_W}
        \(c_t^{H\to W}\)：
        総消費指数\(c_t^{H\to W}\)の図
        \ref{fig:chapter5_beta_shock_graphs_c_H_W}
        を見ると、POLTの総消費指数\(c_t^{H\to W}\)は2つの政策群の中間的な軌道を描いていることがわかる。
        100期までの前半において、POLTの落ち込みは
        総消費水準目標（CLT）や名目総消費水準目標（NCLT）のようには小さく抑えられていないが、
        消費者物価水準目標（CPLT）や消費者物価インフレ目標（IT-CPI）に見られるような極端な大暴落は免れている。
        50期前後におけるオーバーシュート（プラス方向への乖離）についても、
        CPLTやIT-CPIほど巨大な山を作ることはなく、CLTやNCLTに比べるとやや遅れてピークを迎えた後、
        100期以降の後半は緩やかに定常状態へと収束している。
    \item
    \label{itm:chapter5_beta_shock_policies_polt_y_H}
        \(y_t^H\)：
        生産\(y_t^H\)の図
        \ref{fig:chapter5_beta_shock_graphs_y_H}
        を見ると、POLTの生産\(y_t^H\)も
        総消費指数\(c_t^{H\to W}\)と極めて類似した中間的な軌道を描いていることがわかる。
        100期までの前半において、POLTの落ち込みは
        総消費水準目標（CLT）や名目総消費水準目標（NCLT）のようには小さく抑えられていないが、
        消費者物価水準目標（CPLT）や消費者物価インフレ目標（IT-CPI）に見られるような極端な大暴落は免れている。
        またオーバーシュートも50期前後に小さな山を作るにとどまり、
        CPLTやIT-CPIのような遅くて巨大な山を作ることはない。
        CLTやNCLTよりも低いピークをやや遅れて迎えた後、緩やかに定常状態へと収束している。
    \item
    \label{itm:chapter5_beta_shock_policies_polt_pi_H}
        \(\pi_t^H\)：
        グロス・インフレ率\(\pi_t^H\)の図\
        ref{fig:chapter5_beta_shock_graphs_pi_H}
        を見ると、POLTの定常状態からの乖離は
        消費者物価インフレ目標（IT-CPI）のような極端な初期の下落や、
        総消費水準目標（CLT）のような大きな変動に比べれば小さく抑えられていることがわかる。
        しかし名目総消費水準目標（NCLT）や生産者物価水準目標（PPLT）が達成している
        ゼロ近傍での極めて平坦な推移と比較すると、わずかながら乖離が継続している。
        図
        \ref{fig:chapter5_beta_shock_graphs_p_H}
        に見られる物価の上昇ドリフトを反映して中盤以降はインフレ率がプラス圏で推移し続けており、
        この継続的な乖離が厚生を押し下げる要因となっている。
    \item
    \label{itm:chapter5_beta_shock_policies_polt_Delta_t_minus_1_H}
        \(\Delta_{t-1}^H\)：
        価格分散\(\Delta_{t-1}^H\)の図
        \ref{fig:chapter5_beta_shock_welfare_Delta_H}
        を見ると、POLTの定常状態からの乖離は、総消費水準目標（CLT）の持続的な増大や、
        消費者物価水準目標（CPLT）および消費者物価インフレ目標（IT-CPI）が形成する
        巨大な山と比較すれば小さく抑えられていることがわかる。
        しかし、名目総消費水準目標（NCLT）や生産者物価水準目標（PPLT）が達成している
        ほぼゼロの推移と比較すると、50期から150期付近にかけてなだらかな山を形成しており、
        明確な乖離が存在している。
        この中盤における価格分散の蓄積が生産効率を悪化させ、POLTの厚生を押し下げる一因となっている。
\end{enumerate}